\documentclass[10pt,a4paper]{article}
\usepackage[utf8]{inputenc}
\usepackage[vietnamese]{babel}
\usepackage{amsmath,amssymb,amsfonts}
\usepackage{geometry}
\geometry{a4paper, left=10mm, right=10mm, top=10mm, bottom=10mm}
\usepackage{tikz}
\usetikzlibrary{shapes,arrows.meta,positioning,calc}
\usepackage[most]{tcolorbox}
\usepackage{array,booktabs,multicol}
\usepackage{fontawesome}
\usepackage{helvet}
\renewcommand{\familydefault}{\sfdefault}
\usepackage{parskip}

\definecolor{primary}{RGB}{112, 48, 160}
\definecolor{teal}{RGB}{0, 128, 128}
\definecolor{orange}{RGB}{230, 108, 44}
\definecolor{green}{RGB}{46, 139, 87}
\definecolor{rose}{RGB}{192, 0, 0}

\pagestyle{empty}

\begin{document}

% --- HEADER BANNER ---
\begin{tcolorbox}[
  enhanced,
  colback=violet!8!white,
  colframe=violet!80!black,
  arc=2.5mm, boxrule=1pt,
  left=3mm, right=3mm, top=2mm, bottom=2mm
]
  \begin{tabular*}{\textwidth}{@{\extracolsep{\fill}} lr}
    {\small\bfseries\color{violet!85!black}\faBook\ SỔ TAY KIẾN THỨC TOÁN 2} & 
    {\small\bfseries\color{teal!85!black}\faGraduationCap\ CHƯƠNG TRÌNH GDPT 2018}
  \end{tabular*}
  \vspace{-1mm}
  \begin{center}
    \rule{0.98\textwidth}{0.6pt}
  \end{center}
  \vspace{-2mm}
  \centering
  {\LARGE\bfseries\color{violet!90!black}\faBookmark\ BÀI 74: ÔN TẬP KIỂM ĐẾM SỐ LIỆU VÀ KHẢ NĂNG XẢY RA CỦA MỘT SỰ KIỆN}\\[1mm]
  {\itshape\footnotesize\color{black!75} Tổng kết thu thập dữ liệu, đọc biểu đồ tranh và dự đoán chắc chắn, có thể, không thể}
\end{tcolorbox}

\vspace{-2mm}

\noindent
\begin{minipage}[t]{0.49\textwidth}\vspace{0pt}
\begin{tcolorbox}[
  enhanced,
  colback=violet!4!white,
  colframe=violet!75!black,
  arc=2mm, boxrule=0.8pt,
  left=2mm, right=2mm, top=1.8mm, bottom=1.8mm,
  title={\footnotesize\bfseries\color{white}\faKey\ 1. Trọng tâm Thống kê mô tả}
]
{\footnotesize
\begin{itemize}\setlength{\itemsep}{1pt}\setlength{\leftskip}{-2.5mm}
  \item \textbf{Quy trình 3 bước xử lí số liệu:}
  $$\text{Thu thập} \longrightarrow \text{Phân loại} \longrightarrow \text{Kiểm đếm}$$
  \item \textbf{Biểu đồ tranh:}\\
  - Biểu diễn số lượng qua các hình vẽ/biểu tượng trực quan.\\
  - Đọc và so sánh: tìm nhóm nhiều nhất, ít nhất, chênh lệch giữa các nhóm và tổng số lượng.
  \item \textbf{Lập bảng kiểm đếm:} Dùng các vạch thẳng $|$ để đếm nhanh, chính xác không sót.
\end{itemize}
}
\end{tcolorbox}
\end{minipage}
\hfill
\begin{minipage}[t]{0.49\textwidth}\vspace{0pt}
\begin{tcolorbox}[
  enhanced,
  colback=teal!4!white,
  colframe=teal!75!black,
  arc=2mm, boxrule=0.8pt,
  left=2mm, right=2mm, top=1.8mm, bottom=1.8mm,
  title={\footnotesize\bfseries\color{white}\faSliders\ 2. Trọng tâm Xác suất thực nghiệm}
]
{\footnotesize
\begin{itemize}\setlength{\itemsep}{1pt}\setlength{\leftskip}{-2.5mm}
  \item Ba mức độ khả năng xảy ra của một sự kiện:
  \item \textbf{"Chắc chắn":} Luôn xảy ra trong mọi lần thử (ví dụ: rút $1$ con bài trong bộ bài toàn màu đỏ chắc chắn được màu đỏ).
  \item \textbf{"Có thể":} Có khả năng xảy ra nhưng cũng có thể không (ví dụ: gieo xúc xắc được mặt 6 chấm).
  \item \textbf{"Không thể":} Hoàn toàn không bao giờ xảy ra (ví dụ: lấy quả bóng vàng trong hộp chỉ có bóng xanh và đỏ).
\end{itemize}
}
\end{tcolorbox}
\end{minipage}

\vspace{-1.5mm}
\begin{tcolorbox}[
  enhanced,
  colback=orange!4!white,
  colframe=orange!80!black,
  arc=2mm, boxrule=0.8pt,
  left=2mm, right=2mm, top=1.5mm, bottom=1.5mm,
  title={\footnotesize\bfseries\color{white}\faSitemap\ 3. Sơ đồ liên kết Thống kê và Xác suất sơ cấp Toán 2}
]
\centering
\begin{tikzpicture}[x=1.65cm, y=0.75cm]
  \tikzset{arr/.style={-{Stealth[scale=0.8]}}, every node/.append style={align=center}}
  \node[rectangle, rounded corners=2pt, draw=violet!80!black, fill=violet!15, font=\scriptsize\bfseries] (tk) at (0,0.6) {THỐNG KÊ\\Thu thập $\to$ Đếm $\to$ Biểu đồ};
  \node[rectangle, rounded corners=2pt, draw=teal!80!black, fill=teal!15, font=\scriptsize\bfseries] (xs) at (0,-0.6) {XÁC SUẤT\\Chắc chắn - Có thể - Không thể};

  \node[rectangle, rounded corners=2pt, draw=green!70!black, fill=green!15, font=\footnotesize\bfseries] (kq) at (3.5,0) {
    ỨNG DỤNG THỰC TẾ:\\
    Quan sát khoa học \&\\
    Đưa ra quyết định hợp lí
  };

  \draw[arr, thick] (tk) -| (kq);
  \draw[arr, thick] (xs) -| (kq);
\end{tikzpicture}
\end{tcolorbox}

\vspace{-1.5mm}
\noindent
\begin{minipage}[t]{0.49\textwidth}\vspace{0pt}
\begin{tcolorbox}[
  enhanced,
  colback=green!4!white,
  colframe=green!60!black,
  arc=2mm, boxrule=0.8pt,
  left=2mm, right=2mm, top=1.8mm, bottom=1.8mm,
  title={\footnotesize\bfseries\color{white}\faCheckCircle\ 4. Bài tập thực hành mẫu}
]
{\footnotesize
\textbf{1. Điền từ thích hợp:}\\
Trong hộp có $4$ bút mực tím và $1$ bút mực đen. Lấy ngẫu nhiên $1$ chiếc bút:\\
- Lấy được bút mực tím: \textbf{Có thể}.\\
- Lấy được bút mực đỏ: \textbf{Không thể}.\\
- Lấy được bút viết chữ: \textbf{Chắc chắn}.\\
\textbf{2. Khai thác biểu đồ:}\\
Biểu đồ số xe đạp bán trong tuần: Thứ Hai $5$ xe, Thứ Ba $3$ xe, Thứ Tư $7$ xe.\\
- Ngày bán nhiều nhất: \textbf{Thứ Tư} ($7$ xe).\\
- Tổng xe bán được: $5 + 3 + 7 = \mathbf{15}$ (xe).
}
\end{tcolorbox}
\end{minipage}
\hfill
\begin{minipage}[t]{0.49\textwidth}\vspace{0pt}
\begin{tcolorbox}[
  enhanced,
  colback=rose!4!white,
  colframe=rose!70!black,
  arc=2mm, boxrule=0.8pt,
  left=2mm, right=2mm, top=1.8mm, bottom=1.8mm,
  title={\footnotesize\bfseries\color{white}\faExclamationTriangle\ 5. Lời dặn dò kĩ năng}
]
{\footnotesize
\begin{itemize}\setlength{\itemsep}{1.5pt}\setlength{\leftskip}{-2.5mm}
  \item Khi chọn từ "Chắc chắn" hay "Có thể", phải đọc thật kĩ các loại đồ vật có trong bài.
  \item Đừng nhầm lẫn giữa \textbf{"sự kiện có thể"} và \textbf{"sự kiện chắc chắn"}. Dù khả năng xảy ra rất cao nhưng nếu vẫn có trường hợp khác thì vẫn chỉ là "Có thể".
  \item Đếm biểu tượng tranh cẩn thận, không đếm vội.
\end{itemize}
}
\end{tcolorbox}
\end{minipage}


\end{document}
